% block: shifted-systems-chapter-heading page:175 type:chapter
\clearpage\section{{Applications to Shifted Linear Systems}}\label{chap:shifted-linear-systems}
\setcounter{footnote}{0}
% block: shifted-systems-abstract page:175 type:ordinary-paragraph
\noindent\textbf{{Abstract}}\quad
{As seen in Sect.\ \ref{sec:computational-physics} (computational physics) and Sect.\ \ref{sec:machine-learning} (machine learning), one needs to solve a set of linear systems of the form}
\[(A+\sigma_k I)\bm x^{(k)}=\bm b,\quad k=1,2,\ldots,m,\]
{where $\sigma_k$ is a scalar, which are called shifted linear systems. When using the LU decomposition in Sect.\ \ref{subsec:lu-decomposition}, $m$ times LU decompositions are required. Similarly, when using stationary iterative methods in Sect.\ \ref{sec:stationary-iterations}, solving $m$ linear systems, i.e., $(L+D+\sigma_k I)\bm z=\bm v$ for all $k$, is required. To solve these linear systems, one can apply a suitable Krylov subspace method to each linear system. On the other hand, if the initial residual vector with respect to the $i$th linear system and the initial residual vector with respect to the $j$th linear system are collinear, i.e., the angle between the two initial residual vectors is $0$ or $\pi$, the generated Krylov subspaces become the same, which is referred to as the shift-invariance property of Krylov subspaces. This means that we can share the information of only one Krylov subspace to solve all the shifted linear systems, leading to efficient computations of shifted linear systems. In this chapter, Krylov subspace methods using the shift-invariance property are derived from the Krylov subspace methods in Chap.\ \ref{chap:krylov-classification}.}
